\documentclass[10pt,a4paper]{amsart}
\usepackage[margin=19mm]{geometry}
\usepackage{amsmath,amssymb,mathtools}
\usepackage[scr=rsfs,cal=euler]{mathalfa}
\usepackage{xfrac}
\usepackage[unicode,hidelinks,bookmarks=false]{hyperref}
\newcommand{\ie}{i.e.\ }
\newcommand{\cf}{cf.\ }
\newcommand{\resp}{respectively\ }
\newcommand{\Subsection}{\subsection}
\providecommand{\nobreakdash}{\nobreak\hbox{-}}
\setlength{\emergencystretch}{2em}
\multlinegap0pt
\usepackage{amssymb}
\usepackage{cancel}
\usepackage{xcolor}
\newtheorem{lemma}{Lemma}
\newcommand{\BSig}{\mathsf{B}\Sigma^0}
\newcommand{\M}{\mathcal{M}}
\newcommand{\RT}{\mathsf{RT}}
\newcommand{\WF}{\mathsf{WF}}
\title{Conservation of Ramsey's theorem for pairs\\ and well-foundedness}
\author{Quentin Le Hou\'erou, Ludovic Levy Patey, Keita Yokoyama}
\date{Source: arXiv:2402.11616v3 (CC BY 4.0)}
\begin{document}
\maketitle

\noindent\emph{Source and licence.} Conservation of Ramsey's theorem for pairs\\ and well-foundedness. Version arXiv:2402.11616v3 (CC BY 4.0), distributed by arXiv under the Creative Commons Attribution 4.0 International licence, http://creativecommons.org/licenses/by/4.0/, as recorded in the arXiv licence metadata for this exact version and shown on the abstract page https://arxiv.org/abs/2402.11616v3. Full licence notice: \texttt{licenses/arxiv-2402.11616-CC-BY-4.0.txt}.\par\smallskip

\noindent\textit{Editorial scope.} Counted unit: the article's lemma lemma (source0.tex lines 522--524). The article's complete original proof of it is delivered with it (source0.tex lines 525--543, one \texttt{proof} environment of 19 lines). No mathematics is written, repaired or re-derived: the statement and the proof are the article's own lines, in the article's own notation, and no retained line is substituted or re-flowed.  No further source-local context is needed: every cross-reference of every retained line resolves inside the two delivered spans.  Nothing was omitted from any retained span, so the package carries no omission record.  No retained line contains a citation, so the wrapper carries no bibliography and no bibliography entry.  No retained line contains a figure environment, an image-inclusion command or drawing code, so no image is delivered and the package declares no asset dependency.  Editorial formatting only, in addition: an AMS \texttt{amsart} preamble that restates the article's own declarations and macros used by the retained lines, namely the environments \texttt{lemma} on separate counters, together with the standard packages those lines need.  The retained environments print in this document as Lemma 1 in order of appearance, whereas the article prints the same items with the numbering of its own full document.  The complete native source of the article as distributed in the arXiv e-print for this exact version is bundled unchanged as \texttt{sources/154757/source0.tex}.
\begin{lemma}\label[lemma]{lem:coh-wf-forcing-wf}
Let $(\sigma, \rho)$ be a condition. For every~$k \in M$, there exists some~$n \in M$ an extension $(\tau, \rho) \leq (\sigma, \rho)$ forcing~$\Phi_e^{G \oplus Y}$ not to be an $\alpha$-decreasing sequence of length more than~$n$ for every~$e < k$.
\end{lemma}

\begin{proof}
By twisting the Turing functionals, we can assume that for every~$e, n \in M$, if $\Phi^\sigma_e(n)\downarrow$, then
\begin{itemize}
    \item $\Phi^\sigma_e(m)\downarrow$ for every~$m < n$
    \item $\Phi^\sigma_e(0), \Phi^\sigma_e(1), \dots, \Phi^\sigma_e(m)$ is a strictly $\alpha$-decreasing sequence
\end{itemize}

Suppose the lemma does not hold. Then in particular, for every~$M$-finite string~$\tau$ compatible with $(\sigma, \rho)$ and every~$n \in M$, there is some~$\tau_1 \succeq \tau$ compatible with~$(\sigma, \rho)$ and some~$e < k$ such that $\Phi^{\tau_1 \oplus Y}(n)\downarrow$.

Build $Y$-recursively an infinite sequence $\tau_0 \preceq \tau_1 \preceq \tau_2 \preceq \dots$ such that for every~$n$,
\begin{itemize}
    \item $\tau_n$ is compatible with $(\sigma, \rho)$
    \item $\Phi^{\tau_n \oplus Y}_e(n)\downarrow$ for some~$e < k$ (which depends on~$n$)
\end{itemize}
Since~$\M \models \BSig_2$, $\M \models \forall k\RT^1_k$,
so there is an $M$-infinite set~$Z \in \M$ and some~$e < k$ such that for every~$n \in Z$, $\Phi^{\tau_n \oplus Y}_e(n)\downarrow$.

Then, letting $f(n) = \Phi^{\tau_n \oplus Y}_e(n)$, we have an $M$-infinite $\alpha$-decreasing sequence, contradicting $\M \models \WF(\alpha)$.
\end{proof}

\end{document}
